\clearpage\fancyhead[L]{\small PALLAS.}\hypertarget{native-359-lower}{}\label{native-359-lower}
\begin{center}\textbf{PALLAS.}\par\vspace{4pt}\hrule\vspace{4pt}\fontsize{10}{12}\selectfont Monatliche Correspondenz. Band XXVI. S. 199\;..\;203. 1812. August.\par\vspace{4pt}\hrule\vspace{4pt}\end{center}

\begin{flushright}Göttingen, 4. August 1812.\end{flushright}
\par\indent Ich habe Ihnen schon vor einiger Zeit gemeldet, dass \emph{Pallas} in diesem Jahr zum ersten Mal wieder aufgefunden wurde. Seither habe ich selbst nur wenige Beobachtungen gemacht. Um die Zeit der Opposition stand der Planet für unseren Standort die ganze Nacht zu hoch, um mit unserem besten Instrument mikrometrische Vergleiche anzustellen. Am Mauerquadranten war er überhaupt nicht zu sehen, weil er nur die zehnte Größe hatte. Deshalb begnügte ich mich damit, ihn um diese Zeit ein einziges Mal mit einem schwächeren Instrument am Kreismikrometer zu beobachten. Die dabei bestimmte Deklination war wenig zuverlässig.

\Needspace{11\baselineskip}\par\indent Hier sind alle meine Beobachtungen:

\begin{center}
\fontsize{10}{12}\selectfont\setlength{\tabcolsep}{3pt}\renewcommand{\arraystretch}{1.15}\begin{tabular}{c|c|c|c}
1812 & Mittlere Zeit & Rektaszension & Nördl. Dekl. \\
\hline
April 9 & $11^{\mathrm u}\;32^{\mathrm m}\;27^{\mathrm s}$ & $269^\circ\;5'\;58''4$ & $16^\circ\;38'\;0''\!:\,:$ \\
Mai 2 & 10 \quad 36 \quad 6 & 268 \quad 42 \quad 5.3 & 21 \quad 4 \quad 31.7 \\
4 & 10 \quad 44 \quad 44 & 268 \quad 31 \quad 38.0 & 21 \quad 25 \quad 4.3 \\
Juni 7 & 11 \quad 13 \quad 51 & 262 \quad 50 \quad 51.1 & 25 \quad 2 \quad 56.6 \\
\end{tabular}
\end{center}

\Needspace{9\baselineskip}\par\indent Dr. \textsc{Olbers} hatte in Bremen die folgenden beiden Beobachtungen gemacht:

\begin{center}
\fontsize{10}{12}\selectfont\setlength{\tabcolsep}{3pt}\renewcommand{\arraystretch}{1.15}\begin{tabular}{c|c|c|c}
1812 & Mittlere Zeit & Rektaszension & Nördl. Dekl. \\
\hline
April 3 & $11^{\mathrm u}\;26^{\mathrm m}\;44^{\mathrm s}$ & $268^\circ\;42'\;23''1$ & $15^\circ\;25'\;4''9$ \\
4 & 12 \quad 2 \quad 33 & 268 \quad 47 \quad 12.5 & 15 \quad 37 \quad 41.1 \\
\end{tabular}
\end{center}

\Needspace{12\baselineskip}\par\indent Aus Paris war er so freundlich, mir die folgenden vier Meridianbeobachtungen von Dr. \textsc{Burkhardt} mitzuteilen:

\begin{center}
\fontsize{10}{12}\selectfont\setlength{\tabcolsep}{3pt}\renewcommand{\arraystretch}{1.15}\begin{tabular}{c|c|c|c}
1812 & Mittlere Zeit & Rektaszension & Nördl. Dekl. \\
\hline
Juni 5 & $12^{\mathrm u}\;35^{\mathrm m}\;57^{\mathrm s}4$ & $263^\circ\;15'\;24''0$ & $24^\circ\;58'\;28''5$ \\
6 & 12 \quad 33 \quad 10.5 & 263 \quad 2 \quad 35. & 25 \quad 0 \quad 30. \\
8 & 12 \quad 21 \quad 37.0 & 262 \quad 37 \quad 6.5 & 25 \quad 3 \quad 57.7 \\
11 & 12 \quad 7 \quad 15.7 & 261 \quad 58 \quad 37.8 & 25 \quad 6 \quad 29.3 \\
\end{tabular}
\end{center}

\Needspace{20\baselineskip}\par\indent Die zweite Beobachtung wurde als weniger genau bezeichnet. Das sind alle Beobachtungen aus diesem Jahr, von denen ich erfahren habe. Herr \textsc{Nicolai} hat sie alle mit den Elementen verglichen und folgende Übereinstimmung gefunden:
\SourceLeaf{360}\begin{center}
\emph{Unterschied.}
\end{center}

\begin{center}\fontsize{10}{12}\selectfont\setlength{\tabcolsep}{6pt}\renewcommand{\arraystretch}{1.15}\begin{tabular}{c|c|c|c}
1812 & Rektasz. & Dekl. & Beobachter\\\hline
April 3 & $+30''3$ & $-\!-10''7$ & Olbers\\
4 & $+33.3$ & $-16.8$ & Olbers\\
9 & $+30.6$ &  & Gauss\\
Mai 2 & $+57.6$ & $-5.4$ & Gauss\\
4 & $+49.4$ & $+6.6$ & Gauss\\
\end{tabular}\end{center}
\begin{center}\fontsize{10}{12}\selectfont\setlength{\tabcolsep}{6pt}\renewcommand{\arraystretch}{1.15}\begin{tabular}{c|c|c|c}
1812 & Rektasz. & Dekl. & Beobachter\\\hline
Juni 5 & $+51''4$ & $+6''1$ & Burckhardt\\
6 & $+59.4$ & $+11.7$ & Burckhardt\\
7 & $+56.7$ & $-34.0$ & Gauss\\
8 & $+59.6$ & $+2.4$ & Burckhardt\\
11 & $+62.0$ & $+8.2$ & Burckhardt\\
\end{tabular}\end{center}

\par\indent Dass die Elemente bei der diesjährigen Opposition um 1 Minute abweichen, schreibe ich hauptsächlich dem Einfluss von Saturn und Mars zu. Bei ihnen wurden die Störungen durch Jupiter so genau wie möglich berücksichtigt, und sie stimmten mit allen sieben ersten Oppositionen sehr genau überein. Herr \textsc{Nicolai} hat die von Saturn verursachten Störungen unter meiner Aufsicht bereits einmal vorläufig berechnet. Eine zweite, genauere Rechnung wird bald entweder dieser außerordentlich geschickte und gewandte Rechner oder ich selbst vornehmen. Über ihr Ergebnis werde ich zu gegebener Zeit ausführlicher berichten.

\par\indent Herr \textsc{Nicolai} hat die Opposition aus den Beobachtungen im Juni wie folgt bestimmt.

\begin{center}\fontsize{10}{12}\selectfont\setlength{\tabcolsep}{0pt}\renewcommand{\arraystretch}{1.15}\begin{tabular}{l}
1812 Juni 10. $3^{\mathrm u}\;2^{\mathrm m}\;48^{\mathrm s}$\\Mittlere Zeit am Göttinger Meridian
\end{tabular}\end{center}
\begin{center}\fontsize{10}{12}\selectfont\setlength{\tabcolsep}{2pt}\renewcommand{\arraystretch}{1.15}\begin{tabular}{lcr}
Wahre Länge & $.\;.\;.\;.\;.\;.\;.\;.\;.\;.\;.\;.\;.\;.\;.$ & $259^\circ\;27'\;51''7$\\
Wahre geozentrische Breite & $.\;.\;.\;.\;.\;.\;.\;.\;.\;.$ & $48\;18\;26.1\;\mathrm{N}.$\\
\end{tabular}\end{center}

\Needspace{18\baselineskip}\par\indent Ich habe inzwischen untersucht, wie genau sich alle acht bisher beobachteten Oppositionen noch in Einklang bringen lassen, wenn man nur die Störungen durch Jupiter berücksichtigt. Dazu müssen an den Elementen die folgenden Korrekturen vorgenommen werden:

\begin{center}\fontsize{9}{11}\selectfont\setlength{\tabcolsep}{2pt}\renewcommand{\arraystretch}{1.15}\begin{tabular}{lcr}
Korrektur der Epoche 1810 & $.\;.\;.\;.\;.\;.\;.\;.\;.\;.\;.$ & $-13''31$\\
der täglichen mittleren Bewegung & $.\;.\;.\;.\;.\;.\;.\;.\;.\;.$ & $-0.01200$\\
der Länge des Perihels & $.\;.\;.\;.\;.\;.\;.\;.\;.\;.\;.\;.$ & $-17''57$\\
des Exzentrizitätswinkels & $.\;.\;.\;.\;.\;.\;.\;.\;.\;.\;.\;.\;.$ & $-4.09$\\
des Logarithmus der großen Halbachse & $.\;.\;.\;.\;.\;.\;,$ & $+0.0000045$\\
der Länge des Knotens & $.\;.\;.\;.\;.\;.\;.\;.\;.\;.\;.\;.\;.\;.$ & $+0''80$\\
der Neigung der Bahn & $.\;.\;.\;.\;.\;.\;.\;.\;.\;.\;.\;.\;.\;.$ & $-1''21$\\
\end{tabular}\end{center}

\Needspace{16\baselineskip}\noindent Mit diesen Korrekturen stimmen alle acht Oppositionen wie folgt überein:

\begin{center}
\emph{Unterschied}
\end{center}

\begin{center}
\fontsize{10}{12}\selectfont\setlength{\tabcolsep}{6pt}\renewcommand{\arraystretch}{1.15}\begin{tabular}{c|c|c}
Opposition & Mittlere Länge & Helioz. Breite \\
\hline
1803 & $+5''8$ & $-2''8$ \\
1804 & $-12.9$ & $-9.7$ \\
1805 & $-15.2$ & $-0.7$ \\
1807 & $+8.3$ & $-5.8$ \\
1808 & $+16.0$ & $+11.7$ \\
1809 & $+13.0$ & $-1.5$ \\
1811 & $+11.7$ & $-9.5$ \\
1812 & $-25.9$ & $+6.1$ \\
\end{tabular}
\end{center}

\par\indent Ich füge nur noch hinzu, dass leichte Abweichungen der hier mitgeteilten Zahlen von den bereits anderswo gedruckten auf eine spätere Berichtigung der Rechnung zurückgehen.

\Needspace{13\baselineskip}\par\indent Herr \textsc{Nicolai} hat die Opposition für das nächste Jahr wie folgt \emph{im Voraus} berechnet:
\SourceLeaf{361}\begin{center}\fontsize{10}{12}\selectfont\setlength{\tabcolsep}{0pt}\renewcommand{\arraystretch}{1.15}\begin{tabular}{l}
Zeit 1813 Aug. 18\quad $8^{\mathrm u}\;50^{\mathrm m}\;5^{\mathrm s}$\\Mittlere Zeit in Göttingen
\end{tabular}\end{center}
\begin{center}\fontsize{10}{12}\selectfont\setlength{\tabcolsep}{2pt}\renewcommand{\arraystretch}{1.15}\begin{tabular}{lcr}
Wahre Länge & $.\;.\;.\;.\;.\;.\;.\;.\;.$ & $325^\circ\;24'\;17''9$\\
Geozentr. Breite & $.\;.\;.\;,\;.\;.\;.\;.$ & $24\;37\;23.8\;\mathrm{N}.$\\
\end{tabular}\end{center}

\par\indent Ich habe auch \emph{Juno} mehrmals beobachtet. Sie war dieses Jahr außerordentlich lichtschwach. In meinem nächsten Brief werde ich ausführlich über die daraus abgeleiteten Ergebnisse berichten und eine von Herrn \textsc{Wachter} berechnete Ephemeride für das nächste Jahr mitteilen.
